\section{About this book}

\subsection{The books}

The Missal contains the prayers and directions for Mass. The Breviary
contains the Divine Office or Hours, the Church's daily prayer of psalms,
hymns, lessons and prayers.

The Mass readings, prayers and Roman stations come from the 1962
\textit{Missale Romanum}. The general rubrics were promulgated by John~XXIII
in 1960. Advent begins with First Vespers of the Sunday nearest November~30
and ends after None, the afternoon Hour, on December~24. Christmas begins
with its First Vespers. If December~24 is a Sunday, the Vigil of the
Nativity takes the place of the Fourth Sunday without its commemoration.

For the English of the traditional Office this book draws on the
Marquess of Bute's \textit{Roman Breviary}. Its calendar and rules come
from an earlier edition. The prayer appendix uses its wording of
the O Antiphons and \textit{Alma Redemptoris Mater}. Scripture is the
Douay--Rheims in Challoner's revision. Psalm and verse numbers follow that
Bible. In citations using modern book names, 1 and 2 Samuel are 1 and 2 Kings,
and 1 Kings is 3 Kings in the Douay--Rheims. The daily readings can be used
at home throughout Advent.

The Roman stations are the churches appointed in the Missal for the
public liturgy at Rome. The local parish calendar establishes what Mass
and services will be celebrated where you are.

\subsection{Fast and abstinence}

For the universal Latin Church in 1962, the 1917 Code of Canon Law,
canon~1252, \S\S2 and 4, assigned fast and abstinence to the Ember Days and
the Vigil of the Nativity, with exceptions for their occurrence on a
Sunday or holy day of obligation. On December~3, 1959, John~XXIII
granted the faithful the option of keeping the Vigil's fast and abstinence
on December~23. The Vigil itself kept its liturgical date. Paul~VI
revised the universal discipline in
\textit{Paenitemini} in 1966.

Under the 1983 Code of Canon Law, canons~1249--1253, verified for the Latin
Church on November~20, 2026, Fridays and Lent remain penitential times. The
universal norm prescribes abstinence on Friday unless a solemnity occurs
and fast and abstinence on Ash Wednesday and Good Friday. The Code gives
bishops' conferences authority to determine observance and to make
substitutions. It assigns no fast to the Advent season or to the Advent
Ember Days and Christmas Vigil by reason of those observances.

In the United States the bishops' 1966 pastoral statement kept Friday as a
day of penance and especially commends freely chosen abstinence from meat on
ordinary Fridays outside Lent. The accompanying
complementary norm keeps the Friday and Lenten directions in effect.
The statement encourages prayer and participation in the Advent liturgy
and commends freely chosen preparation on the eve of feasts. The advice
on prayer, fasting and alms above is for a reader choosing a personal rule.
